\chapter{Zbiory i operacje na zbiorach}\label{ch:sets}

\section{Dlaczego po logice przechodzimy do zbiorów?}
W rozdziałach~\ref{ch:propositional}--\ref{ch:quantifiers}
nauczyliśmy się odczytywać implikacje, równoważności
i kwantyfikatory. Teraz zastosujemy te narzędzia do opisu
kolekcji obiektów. W matematyce często interesuje nas
najpierw pewna kolekcja obiektów,
a dopiero potem struktury i działania na tej kolekcji. Zbiór punktów
spełniających równanie, zbiór rozwiązań nierówności, rodzina otwartych
podzbiorów przestrzeni czy zbiór wszystkich funkcji o określonych
własnościach to przejawy tej samej idei. Nawet do wypowiedzenia
stwierdzenia „każda liczba z $A$ należy do $B$” potrzebujemy pojęcia
przynależności i zawierania.

\begin{uwaga}[Dwa poziomy ścisłości]\label{rem:sets-metatheory}
W rozdziałach początkowych pracujemy z ustalonymi zbiorami i operacjami,
których istnienie przyjmujemy zgodnie ze zwykłą praktyką matematyczną.
Nie przyjmujemy zasady „każdy warunek $P(x)$ wyznacza zbiór
$\{x:P(x)\}$” bez określenia, skąd pochodzą elementy $x$. Podstawy
aksjomatyczne wrócą wraz z teorią ZFC. Wszystkie poniższe definicje
zbiorów przez warunek należy rozumieć jako wycinanie podzbioru z
wcześniej ustalonego zbioru albo jako skrót dla konstrukcji, której
istnienie jest już uzasadnione.
\end{uwaga}

\begin{definicja}[Przynależność i podzbiór]\label{def:subset}
Niech $A,B$ będą zbiorami. Piszemy $x\in A$, gdy $x$ jest elementem
$A$, oraz $x\notin A$ w przeciwnym przypadku. Mówimy, że $A$ jest
\emph{podzbiorem} $B$, i piszemy $A\subseteq B$, gdy
\[
\forall x\;(x\in A\Longrightarrow x\in B).
\]
Piszemy $A\subsetneq B$, gdy $A\subseteq B$ i $A\ne B$.
\end{definicja}

\begin{przyklad}\label{ex:subset-basic}
Dla $A=\{1,2\}$ i $B=\{1,2,3\}$ mamy $A\subsetneq B$,
ale $\{1,2\}\notin B$: elementy $B$ są liczbami, a nie zbiorem $\{1,2\}$.
Z kolei dla $C=\{\{1,2\},3\}$ zachodzi $\{1,2\}\in C$,
natomiast $1\notin C$. Rozróżnienie $\in$ i $\subseteq$ będzie
niezbędne przy definiowaniu relacji i funkcji.
\end{przyklad}

\begin{definicja}[Zbiór pusty, równość i ekstensjonalność]
Zbiór bez elementów oznaczamy $\varnothing$. Zbiory $A,B$ są równe
wtedy i tylko wtedy, gdy mają dokładnie te same elementy:
\[
A=B\quad\Longleftrightarrow\quad
\forall x\;(x\in A\Longleftrightarrow x\in B).
\]
Ostatnią zasadę nazywamy \emph{ekstensjonalnością zbiorów}.
\end{definicja}

\begin{uwaga}[Co stwierdza ekstensjonalność?]\label{rem:equality-extensionality-xlii}
Znak $=$ jest symbolem równości: oznacza identyczność obiektów.
Kierunek $A=B\Rightarrow\forall x\,(x\in A\Leftrightarrow x\in B)$
wynika z możliwości zastępowania równych obiektów w każdym zdaniu.
Kierunek odwrotny jest zasadą \emph{ekstensjonalności zbiorów}:
zbiory mające te same elementy są równe. Nie należy traktować
całej tej równoważności jako definicji znaku równości w logice;
w aksjomatycznej teorii ZF kierunek odwrotny będzie osobnym
aksjomatem. W niniejszym rozdziale przyjmujemy go jawnie.
\end{uwaga}

\begin{twierdzenie}[Kryterium dwóch zawierań]\label{thm:two-inclusions}
Dla dowolnych zbiorów $A,B$ zachodzi
$A=B$ wtedy i tylko wtedy, gdy $A\subseteq B$ oraz $B\subseteq A$.
Ponadto $\varnothing\subseteq A$ dla każdego zbioru $A$.
\end{twierdzenie}
\begin{proof}
Załóżmy najpierw $A=B$. Każdy element $A$ jest wtedy elementem
$B$, więc $A\subseteq B$; zamieniając role zbiorów otrzymujemy
$B\subseteq A$. Odwrotnie, załóżmy oba zawierania i ustalmy
dowolny obiekt $x$. Jeżeli $x\in A$, pierwsze zawieranie daje
$x\in B$. Jeżeli $x\in B$, drugie daje $x\in A$. Zatem
$x\in A\Longleftrightarrow x\in B$ dla każdego $x$, skąd $A=B$
na mocy ekstensjonalności. Wreszcie warunek
$x\in\varnothing\Longrightarrow x\in A$ jest spełniony dla
każdego $x$: nie istnieje $x\in\varnothing$, który mógłby go
naruszyć. To dowodzi ostatniej tezy.
\end{proof}

\begin{uwaga}[O kwantyfikatorze i pustym zbiorze]
W zdaniu „każdy element zbioru pustego ma własność $P$” nie jest
ukryte twierdzenie o własnościach jakiegoś szczególnego elementu.
Zdanie jest prawdziwe, ponieważ nie istnieje kontrprzykład.
Natomiast „istnieje element zbioru pustego o własności $P$” jest
fałszywe. To dokładnie zachowanie kwantyfikatorów nad pustą dziedziną
omówione w rozdziale~\ref{ch:quantifiers}.
\end{uwaga}


\begin{przyklad}[Czytanie zapisu zbioru przez warunek]
Zapis $A=\{n\in\N:n<3\}$ czytamy: „zbiór tych liczb
naturalnych $n$, dla których $n<3$”. Dla przyjętej konwencji
$\N=\{0,1,2,\ldots\}$ mamy $A=\{0,1,2\}$.
Dwukropek oznacza „takich, że”, a nie dzielenie.
Aby wykazać $n\in A$, musimy sprawdzić dwa fakty:
$n\in\N$ oraz $n<3$. Aby obalić $n\in A$, wystarczy
pokazać, że przynajmniej jeden z tych dwóch warunków zawodzi.
\end{przyklad}

\section{Operacje skończone}\label{sec:set-basic-ops}
\begin{definicja}[Suma, przecięcie, różnica]
Dla zbiorów $A,B$ określamy
\begin{align*}
A\cup B&=\{x:x\in A\text{ lub }x\in B\},\\
A\cap B&=\{x:x\in A\text{ i }x\in B\},\\
A\setminus B&=\{x\in A:x\notin B\}.
\end{align*}
Zbiory $A$ i $B$ nazywamy \emph{rozłącznymi}, gdy
$A\cap B=\varnothing$.
\end{definicja}

\begin{uwaga}[Istnienie zbiorów a zapis warunkowy]\label{rem:finite-set-existence-xlii}
W zapisie $\{x:x\in A\lor x\in B\}$ zmienna $x$ nie
przebiega przez „zbiór wszystkich obiektów”. Jest to skrót
\emph{charakteryzacji} sumy $A\cup B$ przez jej elementy.
W ZF istnienie $A\cup B$ otrzymujemy z aksjomatu pary
i aksjomatu sumy, stosując sumę elementów $\{A,B\}$.
Przecięcie $A\cap B$ i różnicę $A\setminus B$ otrzymujemy
przez wycinanie odpowiednio z $A$ zbiorów
$\{x\in A:x\in B\}$ oraz $\{x\in A:x\notin B\}$.
Nie jest tu potrzebna nieograniczona zasada tworzenia zbioru
z dowolnego warunku. Aksjomaty te omówimy dokładnie
w rozdziale~\ref{ch:zf-axioms}.
\end{uwaga}

\begin{przyklad}[Od zbioru rozwiązań do działań na zbiorach]
Niech $A=\{x\in\R:x^2<4\}=(-2,2)$ oraz
$B=\{x\in\R:x\geq 0\}=[0,\infty)$.
Wtedy $A\cap B=[0,2)$, $A\cup B=(-2,\infty)$,
$A\setminus B=(-2,0)$ oraz $B\setminus A=[2,\infty)$.
Punkt $2$ nie należy do $A$, więc należy do $B\setminus A$.
\end{przyklad}

\begin{center}
\begin{tikzpicture}[scale=.92]
  \draw[thick] (-.85,0) circle (1.25);
  \draw[thick] (.85,0) circle (1.25);
  \begin{scope}
    \clip (-.85,0) circle (1.25);
    \fill[blue!22] (.85,0) circle (1.25);
  \end{scope}
  \draw[thick] (-.85,0) circle (1.25);
  \draw[thick] (.85,0) circle (1.25);
  \node at (-1.47,0) {$A$};
  \node at (1.47,0) {$B$};
  \node[font=\small] at (0,-1.6) {Zacieniowana część to $A\cap B$.};
\end{tikzpicture}
\end{center}
\begin{twierdzenie}[Prawa przemienności, łączności i rozdzielności]
\label{thm:set-lattice}
Dla dowolnych zbiorów $A,B,C$ zachodzą
\begin{align*}
A\cup B&=B\cup A,& A\cap B&=B\cap A,\\
(A\cup B)\cup C&=A\cup(B\cup C),&
(A\cap B)\cap C&=A\cap(B\cap C),\\
A\cap(B\cup C)&=(A\cap B)\cup(A\cap C),&
A\cup(B\cap C)&=(A\cup B)\cap(A\cup C).
\end{align*}
\end{twierdzenie}
\begin{proof}
We wszystkich przypadkach używamy kryterium dwóch zawierań
(Twierdzenie~\ref{thm:two-inclusions}) lub, równoważnie, sprawdzamy
przynależność dowolnego elementu. Dla przemienności sumy:
$x\in A\cup B$ oznacza, że $x\in A$ lub $x\in B$; kolejność
tych dwóch możliwości nie ma znaczenia, więc jest to równoważne
$x\in B\cup A$. Identyczny argument ze spójnikiem „i” daje
przemienność przecięcia.

Dla łączności sumy warunek $x\in (A\cup B)\cup C$
oznacza $(x\in A\lor x\in B)\lor x\in C$, czyli
$x\in A\lor(x\in B\lor x\in C)$, czyli
$x\in A\cup(B\cup C)$. Analogiczna zamiana nawiasów dla
spójnika „i” dowodzi łączności przecięcia.

Wykażmy pierwszą rozdzielność w obu kierunkach. Jeżeli
$x\in A\cap(B\cup C)$, to $x\in A$ oraz
$(x\in B\lor x\in C)$. Jeżeli zachodzi $x\in B$,
to $x\in A\cap B$, w przeciwnym rozpatrywanym przypadku
$x\in C$, więc $x\in A\cap C$. W obu przypadkach
$x\in(A\cap B)\cup(A\cap C)$. Odwrotnie, jeśli $x$ należy
do tej ostatniej sumy, to należy do $A\cap B$ lub do $A\cap C$;
w obu przypadkach $x\in A$ i co najmniej jeden z warunków
$x\in B$, $x\in C$ jest spełniony. Stąd
$x\in A\cap(B\cup C)$. Druga rozdzielność wynika analogicznie:
$x\in A\cup(B\cap C)$ zachodzi wtedy i tylko wtedy, gdy
$(x\in A\lor x\in B)$ oraz $(x\in A\lor x\in C)$,
czyli gdy $x\in(A\cup B)\cap(A\cup C)$.
\end{proof}

\begin{definicja}[Dopełnienie względem zbioru odniesienia]
Jeżeli $A\subseteq U$ i $U$ jest ustalonym zbiorem odniesienia,
to \emph{dopełnieniem $A$ w $U$} nazywamy
$A^{c,U}:=U\setminus A$. Jeśli $U$ jest jasny z kontekstu,
piszemy $A^c$. Dopełnienie zawsze zależy od wyboru $U$.
\end{definicja}

\begin{twierdzenie}[Prawa de Morgana]\label{thm:de-morgan-finite}
Jeżeli $A,B\subseteq U$, to
\[
(A\cup B)^{c,U}=A^{c,U}\cap B^{c,U},\qquad
(A\cap B)^{c,U}=A^{c,U}\cup B^{c,U}.
\]
\end{twierdzenie}
\begin{proof}
Weźmy dowolne $x\in U$. Przynależność
$x\in (A\cup B)^{c,U}$ oznacza, że $x\notin A\cup B$.
Definicja sumy mówi, że nie zachodzi ani $x\in A$, ani
$x\in B$. Ponieważ $x\in U$, są to równoważnie warunki
$x\in A^{c,U}$ i $x\in B^{c,U}$. Otrzymujemy
$x\in A^{c,U}\cap B^{c,U}$. Wszystkie kroki są
równoważnościami, więc udowodniona została równość zbiorów.
W drugim przypadku $x\notin A\cap B$ oznacza,
że $x\notin A$ lub $x\notin B$, co przy $x\in U$
jest równoważne przynależności do sumy dopełnień.
Żaden z porównywanych zbiorów nie ma elementów poza $U$,
więc dowody dla $x\in U$ są wystarczające.
\end{proof}

\begin{twierdzenie}[Różnica zbiorów a dwa spójniki]\label{thm:difference-laws-xlii}
Dla dowolnych zbiorów $A,B,C$ zachodzą tożsamości
\[
(A\setminus B)\setminus C=A\setminus(B\cup C),
\qquad
A\setminus(B\setminus C)=(A\setminus B)\cup(A\cap C).
\]
W szczególności różnica zbiorów nie jest na ogół łączna.
\end{twierdzenie}
\begin{proof}
Ustalmy obiekt $x$. Na mocy definicji różnicy:
\begin{align*}
x\in(A\setminus B)\setminus C
&\iff x\in A\setminus B\text{ i }x\notin C\\
&\iff x\in A\text{ i }x\notin B\text{ i }x\notin C\\
&\iff x\in A\text{ i }x\notin B\cup C\\
&\iff x\in A\setminus(B\cup C).
\end{align*}
Każde przejście jest równoważnością, więc pierwszy wzór
wynika z ekstensjonalności. W drugim wzorze kluczowe
jest zanegowanie koniunkcji określającej $B\setminus C$:
\begin{align*}
x\in A\setminus(B\setminus C)
&\iff (x\in A)\land\neg((x\in B)\land(x\notin C))\\
&\iff (x\in A)\land((x\notin B)\lor(x\in C))\\
&\iff ((x\in A)\land(x\notin B))
  \lor((x\in A)\land(x\in C))\\
&\iff x\in(A\setminus B)\cup(A\cap C).
\end{align*}
W trzecim przejściu użyliśmy rozdzielności koniunkcji
względem alternatywy. Ostatni warunek oznacza dokładnie
przynależność do zbioru po prawej stronie.
\end{proof}
\begin{przyklad}[Kiedy zmiana nawiasów daje inny wynik?]\label{ex:difference-parentheses-xlii}
Dla $A=B=C=\{1\}$ mamy $B\setminus C=\varnothing$,
więc $A\setminus(B\setminus C)=\{1\}$.
Tymczasem $A\setminus B=\varnothing$, a zatem
$(A\setminus B)\setminus C=\varnothing$.
Pierwsza tożsamość twierdzenia oblicza drugie z tych wyrażeń
jako $A\setminus(B\cup C)$; druga tożsamość oblicza
pierwsze jako $(A\setminus B)\cup(A\cap C)$.
\end{przyklad}


\begin{twierdzenie}[Zawieranie a działania na zbiorach]\label{thm:subset-test}
Dla dowolnych zbiorów $A,B$ następujące warunki są równoważne:
\begin{enumerate}[label=(\roman*)]
\item $A\subseteq B$;
\item $A\cap B=A$;
\item $A\cup B=B$.
\end{enumerate}
\end{twierdzenie}
\begin{proof}
Załóżmy $A\subseteq B$. Jeśli $x\in A\cap B$, to
$x\in A$, stąd $A\cap B\subseteq A$. Z drugiej strony
jeśli $x\in A$, to również $x\in B$, więc
$x\in A\cap B$. Otrzymujemy (ii).

Niech teraz zachodzi (ii). Dla dowolnego $x\in A$ mamy
$x\in A\cap B$ z równości zbiorów, więc $x\in B$.
To dowodzi (i).

Następnie z (i) wynika (iii): jeśli $x\in A\cup B$,
to $x\in A$ albo $x\in B$; w pierwszym przypadku
$x\in B$ przez zawieranie, a w drugim bezpośrednio.
Zatem $A\cup B\subseteq B$. Odwrotnie każdy $x\in B$
należy do $A\cup B$, więc mamy równość.
Wreszcie jeśli zachodzi (iii) i $x\in A$, to
$x\in A\cup B=B$, a więc $x\in B$, co daje (i).
\end{proof}

\begin{przyklad}[Szybkie testowanie zawierania]
Dla $A=\{1,3\}$ i $B=\{1,2,3\}$ mamy
$A\cap B=A$ oraz $A\cup B=B$. Jeśli natomiast
$A=\{1,4\}$, to $A\cup B\ne B$, gdyż $4$ występuje
po lewej, lecz nie po prawej.
\end{przyklad}

\begin{definicja}[Różnica symetryczna]
Różnicą symetryczną $A$ i $B$ nazywamy
$A\symdiff B:=(A\setminus B)\cup(B\setminus A)$.
\end{definicja}
\begin{przyklad}
Dla $A=\{1,2,3\}$ i $B=\{2,3,4\}$ mamy
$A\symdiff B=\{1,4\}$. Różnica symetryczna zatrzymuje
dokładnie te elementy, które należą do \emph{jednego} ze zbiorów,
ale nie do obydwu.
\end{przyklad}

\begin{twierdzenie}[Własności różnicy symetrycznej]\label{thm:symdiff}
Dla dowolnych zbiorów $A,B,C$ zachodzi
$A\symdiff B=B\symdiff A$,
$(A\symdiff B)\symdiff C=A\symdiff(B\symdiff C)$,
$A\symdiff A=\varnothing$ i $A\symdiff\varnothing=A$.
\end{twierdzenie}
\begin{proof}
Definicja bezpośrednio pokazuje przemienność. Ponadto żaden
$x$ nie może jednocześnie należeć i nie należeć do $A$,
więc $A\symdiff A=\varnothing$; przy $B=\varnothing$
definicja daje $A\symdiff\varnothing=A$.

Pozostaje łączność. Ustalmy dowolny obiekt $x$.
Niech $p,q,r$ oznaczają odpowiednio zdania
$x\in A$, $x\in B$, $x\in C$. Z definicji różnicy
symetrycznej wynika, że $x\in A\symdiff B$ wtedy
i tylko wtedy, gdy dokładnie jedno spośród $p,q$
jest prawdziwe. Wszystkie osiem możliwych ustawień
wartości logicznych przedstawia tablica:
\[
\begin{array}{ccc|cc}
p&q&r&(p\mathbin{\triangle}q)\mathbin{\triangle}r
  &p\mathbin{\triangle}(q\mathbin{\triangle}r)\\\hline
0&0&0&0&0\\
0&0&1&1&1\\
0&1&0&1&1\\
0&1&1&0&0\\
1&0&0&1&1\\
1&0&1&0&0\\
1&1&0&0&0\\
1&1&1&1&1
\end{array}
\]
W dwóch ostatnich kolumnach symbol $\triangle$ oznacza
operację „dokładnie jedna wartość jest równa $1$”
wykonywaną na wartościach $0,1$, a nie na zbiorach.
Dla przykładu w wierszu $(1,0,1)$ mamy najpierw
$1\triangle0=1$, potem $1\triangle1=0$;
przy przeciwnym nawiasowaniu $0\triangle1=1$,
a następnie $1\triangle1=0$. W każdym wierszu
obie kolumny są zgodne. Na mocy kryterium równości
zbiorów oba nawiasowania mają więc dokładnie
te same elementy.
\end{proof}

\section{Zbiór potęgowy i rodziny indeksowane}\label{sec:families}
\begin{definicja}[Zbiór potęgowy]\label{def:powerset}
Dla zbioru $X$ oznaczamy przez
$\calP(X):=\{A:A\subseteq X\}$ zbiór wszystkich jego podzbiorów.
\end{definicja}
\begin{przyklad}
Dla $X=\{a,b\}$ mamy
$\calP(X)=\{\varnothing,\{a\},\{b\},\{a,b\}\}$.
W szczególności $\varnothing\in\calP(X)$ oraz
$X\in\calP(X)$, także gdy $X=\varnothing$.
Zauważmy, że $\calP(\varnothing)=\{\varnothing\}$ nie jest
zbiorem pustym.
\end{przyklad}

\begin{uwaga}[Element zbioru potęgowego nie musi być elementem $X$]
\label{rem:powerset-membership-xlii}
Niech $X=\{\varnothing\}$. Wówczas $X$ ma jeden element,
a jego zbiór potęgowy to
$\calP(X)=\{\varnothing,\{\varnothing\}\}$.
Zachodzi $\varnothing\in X$ oraz $X\in\calP(X)$,
lecz $X\notin X$: jedynym elementem $X$ jest
$\varnothing$, a $X\ne\varnothing$. Ponadto
$\calP(\varnothing)=\{\varnothing\}$, podczas gdy
$\calP(X)$ ma dwa elementy. Dzięki temu widać różnicę
między relacjami $\in$ i $\subseteq$ także w przypadku
zbiorów, których elementy same są zbiorami.
\end{uwaga}


\begin{twierdzenie}[Liczba podzbiorów zbioru skończonego]
\label{thm:powerset-finite}
Jeżeli zbiór $X$ ma dokładnie $n$ elementów, gdzie $n\in\N$,
to $\calP(X)$ ma dokładnie $2^n$ elementów.
\end{twierdzenie}
\begin{proof}
Przeprowadzimy indukcję po $n$. Jeżeli $n=0$, to $X=\varnothing$,
a $\calP(X)=\{\varnothing\}$ jest jednoelementowy;
$2^0=1$.
Załóżmy, że dla każdego zbioru mającego $n$ elementów
jego zbiór potęgowy ma $2^n$ elementów. Niech $X$ ma $n+1$
elementów. Wybierzmy jeden z nich, oznaczmy go $a$,
i połóżmy $Y=X\setminus\{a\}$. Wtedy $Y$ ma $n$ elementów.
Każdy podzbiór $C\subseteq X$ należy do dokładnie jednej
z dwóch grup: albo $a\notin C$, albo $a\in C$.
W pierwszej grupie znajdują się dokładnie podzbiory $Y$;
na mocy założenia indukcyjnego jest ich $2^n$.
W drugiej grupie każdy zbiór ma jednoznaczną postać
$D\cup\{a\}$, gdzie $D\subseteq Y$. Każdy wybór $D$
daje dokładnie jeden zbiór z drugiej grupy,
a każdy zbiór z tej grupy wyznacza dokładnie jeden
$D$ przez usunięcie $a$. Dlatego i w tej grupie jest
$2^n$ podzbiorów.
Grupy są rozłączne i wyczerpują wszystkie podzbiory $X$.
Łącznie jest ich $2^n+2^n=2^{n+1}$, co zamyka indukcję.
\end{proof}

\begin{twierdzenie}[Zbiór potęgowy i przecięcia]
\label{thm:powerset-intersection}
Dla dowolnych zbiorów $A,B$ zachodzi
$\calP(A\cap B)=\calP(A)\cap\calP(B)$.
Ponadto $\calP(A)\cup\calP(B)\subseteq\calP(A\cup B)$,
lecz drugie zawieranie nie musi być równością.
\end{twierdzenie}
\begin{proof}
Dla dowolnego zbioru $C$ zachodzi
$C\in\calP(A\cap B)$ wtedy i tylko wtedy, gdy
$C\subseteq A\cap B$. To oznacza, że każdy element $C$
należy jednocześnie do $A$ i $B$, czyli
$C\subseteq A$ oraz $C\subseteq B$.
Z definicji zbioru potęgowego jest to równoważne
$C\in\calP(A)\cap\calP(B)$; pierwsza równość jest dowiedziona.
Jeśli zaś $C\in\calP(A)\cup\calP(B)$, to
$C\subseteq A$ lub $C\subseteq B$, a więc zawsze
$C\subseteq A\cup B$. Stąd drugie zawieranie.
Że może ono być ostre, widzimy dla $A=\{1\}$,
$B=\{2\}$: zbiór $\{1,2\}$ należy do
$\calP(A\cup B)$, ale nie do
$\calP(A)\cup\calP(B)$.
\end{proof}

\begin{twierdzenie}[Monotoniczność zbioru potęgowego]
\label{thm:power-monotonic}
Dla dowolnych zbiorów $A,B$ zachodzi
$A\subseteq B$ wtedy i tylko wtedy, gdy
$\calP(A)\subseteq\calP(B)$.
\end{twierdzenie}
\begin{proof}
Jeżeli $A\subseteq B$ i $C\in\calP(A)$,
to $C\subseteq A$, a z przechodniości zawierania
$C\subseteq B$. Stąd $C\in\calP(B)$, a więc
$\calP(A)\subseteq\calP(B)$.
Odwrotnie, niech $\calP(A)\subseteq\calP(B)$.
Ponieważ $A\subseteq A$, mamy $A\in\calP(A)$.
Zatem $A\in\calP(B)$, co z definicji oznacza
$A\subseteq B$.
\end{proof}

\begin{definicja}[Rodzina indeksowana]
Niech $I$ będzie zbiorem i niech dla każdego $i\in I$
ustalono zbiór $A_i$. Kolekcję $(A_i)_{i\in I}$ nazywamy
\emph{rodziną indeksowaną zbiorów}. Ustalamy jeden
zbiór odniesienia $U$ i zakładamy $A_i\subseteq U$
dla wszystkich $i\in I$. Dla dowolnego $I$ definiujemy
\[
\bigcup_{i\in I}A_i
=\{x\in U:\exists i\in I\;x\in A_i\},\qquad
\bigcap_{i\in I}A_i
=\{x\in U:\forall i\in I\;x\in A_i\}.
\]
Przy $I=\varnothing$ powyższe definicje dają
$\bigcup_{i\in\varnothing}A_i=\varnothing$ oraz
$\bigcap_{i\in\varnothing}A_i=U$.
\end{definicja}

\begin{przyklad}[Powtórzenia w rodzinie indeksowanej]
\label{ex:indexed-repetitions-xlii}
Dla $I=\{1,2,3\}$ i $U=\{a,b\}$ określmy
$A_1=\{a\}$, $A_2=\{a\}$, $A_3=\{b\}$.
Indeksy $1$ i $2$ są różne, lecz wyznaczają ten sam zbiór.
Mimo powtórzenia $\bigcup_{i\in I}A_i=\{a,b\}$ oraz
$\bigcap_{i\in I}A_i=\varnothing$. Rodzina indeksowana
nie jest tym samym co zbiór jej różnych wartości:
informacja, który indeks odpowiada któremu zbiorowi,
może mieć znaczenie. Nie potrzebujemy do tego żadnego
wybierania elementów z poszczególnych $A_i$.
\end{przyklad}


\begin{uwaga}[Dlaczego puste przecięcie wymaga otoczenia?]
Dla pustego zbioru indeksów warunek „$x\in A_i$ dla każdego $i$”
jest spełniony niezależnie od $x$. Bez ustalenia, spośród jakich
obiektów wybieramy $x$, nie ma tu jednego zbioru wszystkich takich
obiektów. Dlatego przyjmujemy puste przecięcie względem
konkretnego $U$, a nie „zbioru wszystkich zbiorów”.
\end{uwaga}

\begin{przyklad}[Nie tylko skończone sumy i przecięcia]
Niech $A_n=(-1/n,1/n)\subseteq\R$ dla $n\ge1$.
Wówczas $A_{n+1}\subseteq A_n$ i
$\bigcap_{n\ge1}A_n=\{0\}$: liczba $0$ należy do każdego
$A_n$, a dla $x\ne0$ wybieramy $n>1/|x|$ i otrzymujemy
$x\notin A_n$. Natomiast
$\bigcup_{n\ge1}A_n=(-1,1)$, bo największy zbiór to $A_1$.
\end{przyklad}

\begin{twierdzenie}[Dowolne prawa de Morgana]\label{thm:arbitrary-de-morgan}
Niech $U$ będzie zbiorem, $I$ dowolnym zbiorem indeksów,
a $A_i\subseteq U$ dla każdego $i\in I$. Wtedy
\[
U\setminus\bigcup_{i\in I}A_i
=\bigcap_{i\in I}(U\setminus A_i),\qquad
U\setminus\bigcap_{i\in I}A_i
=\bigcup_{i\in I}(U\setminus A_i),
\]
gdzie przy pustym $I$ stosujemy umowę o przecięciu względem $U$.
\end{twierdzenie}
\begin{proof}
Ustalmy dowolne $x\in U$. Wtedy
\begin{align*}
x\in U\setminus\bigcup_{i\in I}A_i
&\iff \neg\bigl(\exists i\in I:\;x\in A_i\bigr)\\
&\iff \forall i\in I:\;x\notin A_i\\
&\iff x\in\bigcap_{i\in I}(U\setminus A_i).
\end{align*}
W drugim wzorze analogicznie:
\begin{align*}
x\in U\setminus\bigcap_{i\in I}A_i
&\iff \neg\bigl(\forall i\in I:\;x\in A_i\bigr)\\
&\iff \exists i\in I:\;x\notin A_i\\
&\iff x\in\bigcup_{i\in I}(U\setminus A_i).
\end{align*}
W drugim kroku użyliśmy praw negowania kwantyfikatorów
udowodnionych w Twierdzeniu~\ref{thm:negation-quantifiers}. Wszystkie
porównywane zbiory są podzbiorami $U$, więc stosujemy ekstensjonalność.
Dla $I=\varnothing$ pierwsza tożsamość daje
$U\setminus\varnothing=U$, druga
$U\setminus U=\varnothing$, zgodnie z przyjętą umową.
\end{proof}

\begin{twierdzenie}[Rozdzielność względem dowolnych rodzin]
\label{thm:arbitrary-distributivity}
Niech $I$ będzie dowolnym zbiorem indeksów, a $B$ oraz
$A_i$ ($i\in I$) niech będą podzbiorami jednego zbioru $U$.
Przyjmujemy, że przecięcie pustej rodziny podzbiorów $U$ jest równe $U$. Wtedy
\[
B\cap\bigcup_{i\in I}A_i
=\bigcup_{i\in I}(B\cap A_i),\qquad
B\cup\bigcap_{i\in I}A_i
=\bigcap_{i\in I}(B\cup A_i).
\]
\end{twierdzenie}
\begin{proof}
Niech $x\in U$. Aby należeć do pierwszego zbioru, $x$
musi należeć do $B$ i do co najmniej jednego $A_i$.
Oznacza to, że dla pewnego $i\in I$ jednocześnie
$x\in B$ oraz $x\in A_i$; jest to dokładnie warunek
przynależności do sumy po prawej stronie.

Dla drugiego wzoru załóżmy najpierw, że
$x\in B\cup\bigcap_i A_i$. Jeśli $x\in B$, to dla każdego
$i$ należy do $B\cup A_i$. Jeśli $x\notin B$, to z założenia
$x$ należy do wszystkich $A_i$, a zatem także do wszystkich
$B\cup A_i$. W obu przypadkach $x$ należy do przecięcia
po prawej stronie.
Odwrotnie, jeżeli $x$ należy do wszystkich $B\cup A_i$
i $x\in B$, to należy do sumy po lewej. Jeśli zaś $x\notin B$,
to z $x\in B\cup A_i$ dla każdego $i$ wynika $x\in A_i$
dla każdego $i$. Stąd $x\in\bigcap_i A_i$, a więc znowu
$x$ należy do sumy po lewej. Oba zawierania dają równość.
\end{proof}

\begin{przyklad}[Pusta rodzina jako sprawdzian hipotez]
Jeżeli $I=\varnothing$, to w pierwszej równości twierdzenia obie
strony są puste: po lewej $B\cap\varnothing$, po prawej suma
pustej rodziny. W drugiej równości po lewej otrzymujemy
$B\cup U=U$, a po prawej przecięcie pustej rodziny, czyli $U$.
Zatem obie równości pozostają prawdziwe bez hipotezy $I\ne\varnothing$.
\end{przyklad}

\begin{twierdzenie}[Różnica i dowolne sumy oraz przecięcia]
\label{thm:relative-complement-families}
Niech $U$ będzie zbiorem, $A\subseteq U$,
$I$ zbiorem indeksów oraz $B_i\subseteq U$ dla każdego $i\in I$.
Wówczas, przy konwencji przecięcia pustej rodziny względem $U$,
\[
 A\setminus\bigcup_{i\in I}B_i=\bigcap_{i\in I}(A\setminus B_i)
 \quad\text{dla }I\ne\varnothing,
 \qquad
 A\setminus\bigcap_{i\in I}B_i=\bigcup_{i\in I}(A\setminus B_i).
\]
Gdy $I=\varnothing$, pierwsza równość zachodzi po zastąpieniu
prawego pustego przecięcia przecięciem \emph{względem $A$},
a druga zachodzi bez zmian.
\end{twierdzenie}
\begin{proof}
Gdy $I\ne\varnothing$, dla dowolnego $x\in U$ mamy
\begin{align*}
 x\in A\setminus\bigcup_i B_i
 &\iff (x\in A)\land\neg\exists i\in I\ (x\in B_i)\\
 &\iff (x\in A)\land\forall i\in I\ (x\notin B_i)\\
 &\iff \forall i\in I\ ((x\in A)\land(x\notin B_i))\\
 &\iff x\in\bigcap_i(A\setminus B_i).
\end{align*}
Przejście z drugiego do trzeciego warunku wymaga $I\ne\varnothing$:
stały warunek $x\in A$ można wówczas przenieść pod kwantyfikator
$\forall i\in I$, a z otrzymanego warunku uniwersalnego
odzyskać $x\in A$. Przejście z trzeciego do czwartego warunku
wynika już z definicji przecięcia względem $U$.
Jeśli $I$ jest puste, lewa strona pierwszej równości
to $A$, a przecięcie puste względem $U$ to $U$; dlatego bez
wskazanej modyfikacji równość może być fałszywa.
Dla drugiej równości nie trzeba wymagać niepustości $I$:
\begin{align*}
 x\in A\setminus\bigcap_i B_i
 &\iff (x\in A)\land\neg\forall i\in I\ (x\in B_i)\\
 &\iff (x\in A)\land\exists i\in I\ (x\notin B_i)\\
 &\iff \exists i\in I\ ((x\in A)\land(x\notin B_i))\\
 &\iff x\in\bigcup_i(A\setminus B_i).
\end{align*}
W przypadku pustego $I$ zdania egzystencjalne są fałszywe,
więc obie strony drugiej równości są puste.
\end{proof}

\begin{uwaga}[Poprawna konwencja przy różnicach i pustej rodzinie]
\label{rem:empty-relative-intersection-xlii}
Dla dowolnego $I$ można zapisać pierwszy wzór
Twierdzenia~\ref{thm:relative-complement-families}
bez wyjątku, jeśli \emph{wszystkie} zbiory $A\setminus B_i$
rozpatrujemy jako podzbiory $A$, a puste przecięcie
w tym wzorze liczymy względem $A$:
\[
A\setminus\bigcup_{i\in I}B_i
=\{x\in A:\forall i\in I\;x\notin B_i\}.
\]
Dla $I=\varnothing$ prawa strona jest całym $A$.
Jeśli natomiast puste przecięcie po prawej stronie
zdefiniowano względem większego $U$, wynosi ono $U$;
wtedy równanie nie jest prawdziwe, gdy $A\ne U$.
Przeniesienie stałego warunku $x\in A$ pod kwantyfikator
$\forall i\in I$ jest równoważnym przekształceniem warunków
dla $x\in U$ tylko przy $I\ne\varnothing$.
Dla pustego $I$ można zachować równość \emph{zbiorów},
licząc przecięcie względem $A$; nie oznacza to jednak
równoważności obu warunków dla dowolnego $x\in U$.
\end{uwaga}


\section{Ciągi zbiorów i granice mnogościowe}\label{sec:set-limits}
W analizie często obserwujemy nie pojedynczy zbiór, ale zbiór
zmieniający się wraz z numerem etapu. Dwa naturalne pytania to:
„które punkty pozostają w zbiorach od pewnego momentu?” i
„które punkty wracają do zbiorów nieskończenie wiele razy?”.

\begin{definicja}[Granica dolna i górna zbiorów]\label{def:set-liminf-limsup}
Niech $(A_n)_{n\ge1}$ będzie ciągiem podzbiorów zbioru $U$.
Definiujemy
\[
\liminf_{n\to\infty}A_n
:=\bigcup_{m\ge1}\bigcap_{n\ge m}A_n,\qquad
\limsup_{n\to\infty}A_n
:=\bigcap_{m\ge1}\bigcup_{n\ge m}A_n.
\]
Pierwszy zbiór zawiera elementy należące do wszystkich $A_n$
od pewnego miejsca; drugi --- elementy należące do nieskończenie
wielu $A_n$.
\end{definicja}

\begin{twierdzenie}[Znaczenie obu granic]\label{thm:set-lims}
Niech $(A_n)_{n\ge1}$ będzie ciągiem podzbiorów $U$. Wtedy:
\begin{enumerate}[label=(\roman*)]
\item $x\in\liminf_n A_n$ wtedy i tylko wtedy, gdy istnieje
$m\ge1$ takie, że $x\in A_n$ dla wszystkich $n\ge m$;
\item $x\in\limsup_n A_n$ wtedy i tylko wtedy, gdy dla każdego
$m\ge1$ istnieje $n\ge m$ z $x\in A_n$;
\item $\liminf_n A_n\subseteq\limsup_n A_n$;
\item $U\setminus\limsup_n A_n
=\liminf_n (U\setminus A_n)$ i
$U\setminus\liminf_n A_n
=\limsup_n(U\setminus A_n)$.
\end{enumerate}
\end{twierdzenie}
\begin{proof}
Dla (i) rozwijamy definicję sumy:
$x\in\bigcup_m\bigcap_{n\ge m}A_n$ oznacza istnienie $m$
takiego, że $x$ należy do całego przecięcia po $n\ge m$;
to z definicji przecięcia znaczy $x\in A_n$ dla wszystkich
$n\ge m$. W (ii) kolejność kwantyfikatorów jest odwrotna:
$x$ musi należeć do $\bigcup_{n\ge m}A_n$ dla każdego
$m$, zatem dla każdego $m$ można wskazać odpowiednie $n\ge m$.

Jeżeli zachodzi (i), wybierzmy odpowiadające mu $m_0$.
Dla dowolnego $m$ weźmy $n\ge\max\{m,m_0\}$.
Wówczas $x\in A_n$, więc spełniony jest warunek (ii).
To pokazuje (iii). Wreszcie (iv) otrzymujemy dwukrotnie
stosując dowolne prawa de Morgana
(Twierdzenie~\ref{thm:arbitrary-de-morgan}); dla przykładu
\begin{align*}
U\setminus\limsup_n A_n
&=U\setminus\bigcap_m\bigcup_{n\ge m}A_n\\
&=\bigcup_m\left(U\setminus\bigcup_{n\ge m}A_n\right)\\
&=\bigcup_m\bigcap_{n\ge m}(U\setminus A_n).
\end{align*}
Ostatni zbiór to $\liminf_n(U\setminus A_n)$.
Dla drugiej tożsamości zapisujemy osobno wszystkie kroki:
\begin{align*}
U\setminus\liminf_n A_n
&=U\setminus\bigcup_m\bigcap_{n\ge m}A_n\\
&=\bigcap_m\left(U\setminus\bigcap_{n\ge m}A_n\right)\\
&=\bigcap_m\bigcup_{n\ge m}(U\setminus A_n)\\
&=\limsup_n(U\setminus A_n).
\end{align*}
W obu dowodach zastosowaliśmy prawo de Morgana do rodziny
zewnętrznej i ponownie do rodziny wewnętrznej.
\end{proof}

\begin{przyklad}[Zmienność bez granicy]
Niech $A_n=\{0\}$ dla parzystych $n$ oraz $A_n=\{1\}$
dla nieparzystych $n$. Ani $0$, ani $1$ nie występują w
$A_n$ od pewnego miejsca zawsze, więc $\liminf_n A_n=\varnothing$.
Obie liczby należą jednak do nieskończenie wielu zbiorów,
więc $\limsup_n A_n=\{0,1\}$.
\end{przyklad}

\begin{przyklad}[Przedziały skupiające się wokół zera]
\label{ex:set-limits-boundary-xlii}
Rozważmy w $U=\R$ dwa ciągi:
$A_n=[-1/n,1/n]$ oraz $B_n=(0,1/n)$ dla $n\ge1$.
W pierwszym przypadku $0\in A_n$ dla każdego $n$.
Jeśli $x\ne0$, wybierzmy liczbę naturalną $N>1/|x|$.
Dla każdego $n\ge N$ zachodzi $1/n<|x|$,
więc $x\notin A_n$. Z definicji granic wynika
$\liminf A_n=\limsup A_n=\{0\}$.
W drugim przypadku $0\notin B_n$ dla żadnego $n$;
liczby niedodatnie również do nich nie należą.
Jeżeli $x>0$, wybieramy $N>1/x$. Wtedy dla wszystkich
$n\ge N$ mamy $1/n<x$, czyli $x\notin B_n$.
Każda liczba rzeczywista należy więc do co najwyżej
skończenie wielu zbiorów $B_n$, stąd
$\liminf B_n=\limsup B_n=\varnothing$.
Przedziały $B_n$ zbliżają się geometrycznie do zera,
lecz granica mnogościowa nie zawiera zera: brak
przynależności w każdym kroku jest tu decydujący.
\end{przyklad}


\begin{twierdzenie}[Monotoniczne ciągi zbiorów]\label{thm:monotone-sets}
Jeżeli $A_n\subseteq A_{n+1}$ dla wszystkich $n\ge1$,
to $\liminf_n A_n=\limsup_n A_n=\bigcup_{n\ge1}A_n$.
Jeżeli natomiast $A_{n+1}\subseteq A_n$ dla wszystkich
$n\ge1$, to $\liminf_n A_n=\limsup_n A_n=\bigcap_{n\ge1}A_n$.
\end{twierdzenie}
\begin{proof}
Załóżmy, że ciąg jest rosnący. Jeśli $x$ należy do
$\bigcup_n A_n$, to istnieje $m$ takie, że $x\in A_m$.
Przez monotoniczność $x\in A_n$ dla wszystkich $n\ge m$,
a zatem $x\in\liminf_n A_n$ na mocy
Twierdzenia~\ref{thm:set-lims}(i). Odwrotnie,
$x\in\liminf_n A_n$ implikuje przynależność do przynajmniej
jednego $A_m$, więc do ich sumy. Mamy równość
$\liminf_n A_n=\bigcup_n A_n$. Ponieważ
$\liminf_n A_n\subseteq\limsup_n A_n$ i każdy element
$\limsup_n A_n$ należy do przynajmniej jednego $A_n$
(na przykład do któregoś $A_n$ z $n\ge1$), otrzymujemy też
równość z granicą górną.

Teraz niech ciąg będzie malejący. Jeśli $x\in\bigcap_n A_n$,
to $x$ należy do każdego $A_n$, więc do obu granic.
Jeżeli $x\notin\bigcap_n A_n$, istnieje $m$ takie, że
$x\notin A_m$. Z malejącości wynika $x\notin A_n$
dla wszystkich $n\ge m$. Wtedy $x\notin\limsup_n A_n$
na mocy Twierdzenia~\ref{thm:set-lims}(ii).
Ponieważ granica dolna zawiera się w górnej, dostajemy
obydwie żądane równości.
\end{proof}

\begin{przyklad}[Przecięcie pustej rodziny a logika]
Niech $U=\{a,b\}$ i $I=\varnothing$. Nie ma ani
jednego zbioru $A_i$ do sprawdzenia. Zdanie
„dla każdego $i\in I$ zachodzi $x\in A_i$”
jest prawdziwe dla $x=a$ i $x=b$, dlatego
$\bigcap_{i\in I}A_i=U$. Natomiast zdanie
„istnieje $i\in I$ z $x\in A_i$” jest dla obu
$x$ fałszywe, więc $\bigcup_{i\in I}A_i=\varnothing$.
To zastosowanie rozdziału~\ref{ch:quantifiers}, a nie
wyjątek od logiki.
\end{przyklad}

\begin{twierdzenie}[Niezmienniczość granic mnogościowych przy skończonej zmianie]\label{thm:lxiv-finite-set-modifications}
Niech $U$ będzie zbiorem, a $(A_n)_{n\ge1}$ i $(B_n)_{n\ge1}$
ciągami jego podzbiorów. Jeżeli istnieje $N\ge1$ takie, że
$A_n=B_n$ dla każdego $n\ge N$, to
$\liminf_n A_n=\liminf_n B_n$ oraz
$\limsup_n A_n=\limsup_n B_n$.
\end{twierdzenie}
\begin{proof}
Ustalmy $x\in U$. Na mocy Twierdzenia~\ref{thm:set-lims}(i)
warunek $x\in\liminf_n A_n$ oznacza, że istnieje $m\ge1$,
dla którego $x\in A_n$ przy każdym $n\ge m$.
Wówczas dla każdego $n\ge\max\{m,N\}$ mamy $x\in B_n$,
więc $x\in\liminf_n B_n$. Zamieniając role obu ciągów,
otrzymujemy także zawieranie odwrotne.

Podobnie, jeżeli $x\in\limsup_n A_n$, to dla dowolnego
$m\ge1$ istnieje $n\ge\max\{m,N\}$ z $x\in A_n$:
wynika to z Twierdzenia~\ref{thm:set-lims}(ii) zastosowanego
do indeksu $\max\{m,N\}$. Ponieważ $n\ge N$, zachodzi
$x\in B_n$. Stąd $x\in\limsup_n B_n$; zamiana ról
ciągów daje drugie zawieranie. Nie korzystaliśmy z aksjomatu wyboru.
\end{proof}
\begin{przyklad}[Zmiana pierwszego wyrazu a granice]\label{ex:lxiv-finite-set-modifications}
Niech $U=\{0,1\}$, $A_1=\{0\}$, $B_1=\{1\}$,
a dla każdego $n\ge2$ niech $A_n=B_n=\{0\}$.
Wtedy $\bigcup_n A_n=\{0\}$, natomiast
$\bigcup_n B_n=\{0,1\}$, więc zwykłe sumy są różne.
Jednak obie granice dolne i obie granice górne są równe $\{0\}$:
element $0$ należy do wszystkich zbiorów od drugiego wyrazu,
a element $1$ występuje tylko raz, w $B_1$.
\end{przyklad}

\section{Dwa ostrzeżenia na zakończenie}
\begin{przyklad}[Pozorne podobieństwo dwóch operacji]
Nie należy traktować odejmowania zbiorów tak jak odejmowania
liczb: znak $\setminus$ nie jest łączny. Pełny kontrprzykład,
a także poprawne wzory z oboma nawiasowaniami,
zapisano w Przykładzie~\ref{ex:difference-parentheses-xlii}
i Twierdzeniu~\ref{thm:difference-laws-xlii}.
\end{przyklad}

\begin{uwaga}[Paradoks Russella]\label{rem:russell}
Przypuśćmy na chwilę, że wolno tworzyć zbiór
$R=\{x:x\notin x\}$ obejmujący wszystkie obiekty o takiej
własności. Pytanie $R\in R$ prowadzi wówczas do równoważności
$R\in R\iff R\notin R$, czyli do sprzeczności.
Nie oznacza to, że działania na konkretnych zbiorach są
sprzeczne. Oznacza, że zakres dopuszczalnych konstrukcji zbiorów
musi zostać określony aksjomatycznie. Później zobaczymy,
że oddzielanie \emph{wewnątrz ustalonego zbioru} jest bezpieczną
operacją dopuszczoną przez ZF.
\end{uwaga}
